\chapter{Conjuntos Abstractos y Espacios}
\label{conjuntos}

Originalmente, Lawvere propuso los siguientes $8$ axiomas para caracterizar a una categoría
de conjuntos abstractos \cite{lawvere1964elementary}. Se conoce a esta axiomática como $\textbf{ETCS}$
o Teoría Elemental de la Categoría de Conjuntos.

\begin{axiom}[Raíces finitas]
Existen todas las raíces (límites y colimites) finitas. Explícitamente:
existen un objeto terminal $1$ y un objeto inicial $0$, para cada par de objetos
$A,B$ existen el producto $A\times B$ y la suma (coproducto) $A+B$, y para cada
par de flechas paralelas $f,g\colon A\to B$ existen su igualador y su coigualador.
\end{axiom}

\begin{axiom}[Exponenciación]
Para cada par de objetos $A,B$ existe el objeto exponencial $B^{A}$.
\end{axiom}

\begin{axiom}[Objeto de Dedekind--Peano]
Existe un objeto $(N,o,s)$ de números naturales
\end{axiom}

\begin{axiom}[$1$ es un separador]
Si $f,g\colon A\to B$ son flechas distintas, entonces existe un elemento
$x\colon 1\to A$ tal que $fx\neq gx$.
\end{axiom}

\begin{axiom}[Axioma de elección]
Si $f\colon A\to B$ es una flecha y $A$ tiene al menos un elemento, entonces
existe una flecha $g\colon B\to A$ tal que $fgf=f$.
\end{axiom}

\begin{axiom}[Los objetos no iniciales tienen elementos]
Si $A$ no es un objeto inicial, entonces existe al menos un elemento
$x\colon 1\to A$.
\end{axiom}

\begin{axiom}[Elementos de una suma]
Sean $i_A\colon A\to A+B$ e $i_B\colon B\to A+B$ las inyecciones canónicas.
Todo elemento $x\colon 1\to A+B$ factoriza a través de $i_A$ o a través de
$i_B$.
\end{axiom}

\begin{axiom}[Existencia de un objeto con más de un elemento]
Existe un objeto que tiene más de un elemento. Equivalentemente,
$1$ y $0$ no son isomorfos.
\end{axiom}

Con el desarrollo de la teoría de topos se demostró que la lista se puede simplificar a pedir
que la categoría sea un topos bien punteado, con objeto de números naturales y axioma de elección.
A partir de esos axiomas, Lawvere desmostró que es válido el axioma de extensionalidad, del par, de la unión
y el axioma de potencia.También demostró un esquema restringido de separación y, por supuesto, la existencia
de un conjunto vacio, el axioma de Elección y el de Infinito, son inmediatos. El axioma de Reemplazo y el
axioma de buena fundación no forman parte de la teoría, aún así, todo modelo de ETCS es
equivalente a $\con$ como categoría \cite{lawvere1964elementary}.

\begin{definition}
  Un objeto de números naturales (ONN) en una categoría con objeto terminal, es una tercia $(N,o,s)$
  donde $N$ es un objeto de la categoría, $o\colon 1\to N$ y $s\colon N\to N$ son morfismos tales que para todo
  objeto $A$ y morfismos $a\colon 1\to A$ y $g\colon A\to A$ existe un único morfismo $f\colon N\to A$ tal que
  el siguiente diagrama conmuta:
  \begin{equation*}
    \begin{tikzcd}
      1\ar{r}{o}\ar[swap]{rd}{a} & N\ar{d}{f}\ar{r}{s} & N\ar{d}{f} \\
      & A\ar[swap]{r}{g} & A.
    \end{tikzcd}
  \end{equation*}
\end{definition}

Intuitivamente, la flecha $o$ de la definición señala al cero y $s$ es la función sucesor, por lo que
$f$ es la función definida recursivamente tal que $f(0)=a$ y $f(n+1)=g(f(n))$.

En una categoría cartesiana cerrada, un ONN permite también hacer recursión con parámetros,
es decir, definir funciones usando que $f(a,0)=g(a)$ y $f(a,n+1)=h(a,f(a,n))$, ver \cite{johnstone2002sketches}.

\begin{proposition}[label=prop:onncoproducto]
  Sea $(N,o,s)$ un ONN en una categoría $\mathcal{C}$ con coproductos binarios, entonces
  \begin{equation*}
    \begin{tikzcd}
      1\ar{r}{o} & N & N\ar[swap]{l}{s}
    \end{tikzcd}
  \end{equation*}
  es un coproducto.
\end{proposition}
\begin{proof}
  La demostración se resume en que el ONN es el objeto inicial de la categoría de T-álgebras con
  $T\colon \mathcal{C}\to \mathcal{C}$ el endofuntor $1+(-)$. Dada una T-álgebra $(A,a)$, con $a\colon 1+A\to A$,
  por definición de coproducto tenemos $i_1\colon 1\to 1+A$ e $i_A\colon A\to
  1+A$, entonces existe una única $f\colon N\to A$ que hace conmutar al siguiente diagrama:
  \begin{equation*}
    \begin{tikzcd}
      1\ar{r}{o}\ar[swap]{rd}{ai_1} & N\ar[dashed]{d}{f}\ar{r}{s} & N\ar[dashed]{d}{f} \\
      & A\ar[swap]{r}{ai_A} & A.
    \end{tikzcd}
  \end{equation*}
  Para que $f$ realmente sea un morfismo entre $(N,[o,s]\colon 1+N\to N)$ y
  $(A,a\colon 1+A\to A)$ necesitamos que $f[o,s]=aT(f)$.
  \pdfcomment{Tal vez el punto es una como y sigue con ``como se muestra en los siguientes diagramas''}
  \begin{equation*}
    \begin{tikzcd}
      1\ar{r}{i_1}\ar[swap]{rdd}{ai_1}\ar{rd}{o} & 1+N\ar{d}{[o,s]} & N\ar[swap]{l}{i_N}\ar{ldd}{fs}\ar{ld}{s} & {} & 1\ar{r}{i_1}\ar[swap]{rdd}{ai_1}\ar{rd}{i_1} & 1+N\ar[swap]{d}{Tf} & N\ar[swap]{l}{i_N}\ar{ldd}{fs}\ar[swap]{ld}{i_Af} \\
        & N\ar{d}{f} & {} & {} & {} & 1+A\ar{d}{a} & {} \\
        & A, & {} & {} & {} & A. & {}
    \end{tikzcd}
  \end{equation*}
  La construcción de $f$ nos permite concluir que los diagramas son conmutativos, así que
  por la propiedad universal del coproducto obtenemos lo que buscábamos y por tanto $(N,[o,s])$ es el
  objeto inicial de las T-álgebras. Pero el morfismo estructural de una T-álgebra inicial es un isomorfismo,
  gracias a que $(TN, T[o,s])$ es también una T-álgebra y no es difícil verificar que el morfismo inducido $g\colon N\to TN$ es
  inverso por ambos lados de $[o,s]$ (en general, del morfismo estructural de
  una T-álgebra inicial, ver \cite{johnstone2002sketches}).
  De esta manera, $N\cong 1+N$.
\end{proof}

\begin{definition}[label=def:predecesor]
  En una categoría con ONN y productos binarios, es posible definir al morfismo predecesor
  como el único $p\colon N\to N$ que hace conmutar al siguiente diagrama:
  \begin{equation*}
    \begin{tikzcd}
      1\ar{r}{o}\ar[swap]{rd}{o\times o} & N\ar[dashed]{d}{(\id,p)}\ar{r}{s} & N\ar[dashed]{d}{(\id,p)} \\
      & N\times N\ar[swap]{r}{(s\pi_1,\pi_1)} & N\times N.
    \end{tikzcd}
  \end{equation*}
\end{definition}

De manera informal, la conmutatividad del diagrama nos dice que $p(0)=0$ y $p(n+1)=n$.

\begin{proposition}[label=prop:sucesormono]
  Sea $(N,o,s)$ un ONN en una categoría con productos binarios, entonces $s$ es un monomorfismo.
\end{proposition}
\begin{proof}
  Sea $p\colon N\to N$ el predecesor de~\ref{def:predecesor}, notemos en la construcción de $p$ que
  la conmutatividad en la segunda entrada dice que $ps=\pi_1(\id_N,p)=\id_N$. Por tanto,
  $s$ tiene retracto y así es mono.
\end{proof}

\begin{proposition}[label=prop:complementosucesor]
  Sea $(N,o,s)$ un ONN en un topos, entonces $o$ es el complemento de $s$ en $\sub(N)$.
\end{proposition}
\begin{proof}
  Ya vimos que $s$ es mono en \ref{prop:sucesormono} y $o$ es trivialmente mono. Por \ref{prop:onncoproducto}
  $N\cong 1+N$, pero en un topos los coproductos son disjuntos, y por tanto $o$ es el complemento de $s$.
\end{proof}

\begin{proposition}
  En toda categoría con ONN podemos hacer ``inducción'', es decir, si $m\colon N'\to N$ es un subobjeto
  tal que tanto $o$ como la restricción de $s$ a $N'$ se factorizan a través de $N'$, entonces $N'\cong N$.
\end{proposition}
\begin{proof}
  Sean $o'\colon 1\to N'$ y $s'\colon N'\to N'$ las factorizaciones, por la propiedad universal del ONN, existe un único 
  morfismo $f\colon N\to N'$ tal que el siguiente diagrama es conmutativo.
  \begin{equation*}
    \begin{tikzcd}
      1\ar{r}{o}\ar{rd}{o'}\ar[swap]{rdd}{o} & N\ar{d}{f}\ar{r}{s} & N\ar{d}{f} \\
      & N'\ar{r}{s'}\ar{d}{m} & N'\ar{d}{m} \\
      & N\ar{r}{s} & N.
    \end{tikzcd}
  \end{equation*}
  Por lo que $mf\colon N\to N'\to N$ tiene que ser la identidad en $N$, pero esto implica que $m$ además de ser mono tiene
  inverso por la derecha, así que $m$ es un isomorfismo.
\end{proof}

\begin{definition}
    Un topos es bien punteado si
    \begin{itemize}
        \item No sucede que $1\cong 0$.
        \item El objeto términal es separador (axioma 4 de ETCS).
    \end{itemize}
\end{definition}

Para enunciar el siguiente resultado, necesitamos más definiciones.

\begin{definition}
    Dado un objeto $A$, el soporte de $A$, $\supp(A)$, es la imagen de la única flecha $A\to 1$, es decir,
    $A\twoheadrightarrow \supp(A)\rightarrowtail 1$ es su factorización epi mono.
\end{definition}

\begin{definition}
    En un topos los soportes se escinden cuando, para todo objeto $A$ el epi $A\to\supp(A)$ tiene una sección.
\end{definition}

Un resultado estándar en teoría de topos caracteriza a los topos bien punteados \cite{maclane2012sheaves}. Resulta que:

\begin{theorem}
    Un topos es bien punteado si y solo si es bivaluado, booleano y los soportes se escinden.
\end{theorem}

Como mencionamos anteriormente, ETCS demuestra un fragmento de los axiomas de ZFC, pero dicho fragmento también prueba a los axiomas de ETCS,
este pedazo de ZFC se le llama BZC (Bounded Zermelo with Choice) o a veces "Teoría de Conjuntos de Mac Lane".

En esta sección seguiremos a Mclarty \cite{mclarty1988defining} y utilizaremos la siguiente lista de axiomas,
a parte de los axiomas de topos, y a esta teoría la llamaremos CS.

\begin{axiom}[No trivial]
    El objeto inicial no tiene elementos.
\end{axiom}

\begin{axiom}[Nullstellensatz]
  \pdfcomment{Todas las ocurrencias de la palabra Nullstellensatz deben estar
  en itálicas. En español las palabras de origen extrangero se escriben en itálicas}
    Cada objeto, o es inicial o bien tiene elementos.
\end{axiom}

\begin{axiom}[Booleano]
    La ley del tercero excluido es internamente válida.
\end{axiom}

\begin{axiom}[Infinito]
    Existe un objeto de números naturales.
\end{axiom}

\begin{proposition}
    Nullstellensatz implica no trivial.
\end{proposition}
\begin{proof}
    El Nullstellensatz implica que, para todo objeto $A$, $A\cong 0$ si y solo si
    no existe ningún morfismo $1\to A$. Por tanto, un objeto inicial no puede tener elementos.
\end{proof}

\begin{proposition}
    Dado un objeto $A$ en un topos, cualquier morfismo $\supp(A)\to A$ hace que el soporte de $A$ se escinda.
\end{proposition}
\begin{proof}
    Sea $A$ un objeto de un topos, $p\colon A\to \supp(A)$ su soporte y $s\colon \supp(A)\to A$ un morfismo. 
    Consideremos al siguiente diagrama que es conmutativo pues acaba en el objeto terminal:
    \begin{equation*}
        \begin{tikzcd}
            \supp(A)\ar[shift left = 0.3ex]{r}{\text{id}}\ar[swap, shift right = 0.3ex]{r}{ps} & \supp(A)\ar[tail]{r}{} & 1
        \end{tikzcd}
    \end{equation*}
    \pdfcomment{Tal vez vale la pena decir que se concluye ps=id.}
\end{proof}

\begin{definition}[label=def:bivaluado]
  Un topos es bivaluado si su clasificador de subobjetos tiene exactamente dos elementos globales, verdadero
  y falso.
\end{definition}

\pdfcomment{Saque las observaciones y los resultados de la definición. Como no
son parte de la definición, no deberían estar dentro del ambiente de definición.
Además, ¿cómo es que el resultado de los soportes hace que Esp sea bivaluado?}
En particular, un topos bivaluado no es trivial y, con el resultado anterior, $\esp$ es bivaluado.

Vale la pena distinguir entre que un topos sea bivaluado y que su lógica interna sea clásica,
ya que, de manera quizá contraintuitiva, tener dos valores de verdad no lo garantiza. Para analizar la distinción tenemos la siguiente definición.

\begin{definition}
  Un topos $\mathcal{E}$ es booleano si $\Omega_{\mathcal{E}} \cong 1+1$. Esto es equivalente a que para
  cada objeto $A$ de $\mathcal{E}$, $\sub(A)$ es una álgebra de Boole, es decir,
  el tercero excluido es válido internamente.
\end{definition}

Ser bivaluado sí implica tercero excluido para fórmulas sin variables libres, pero esto no basta
ni siquiera para que las leyes clásicas de cuantificadores se cumplan. Por
ejemplo, en $\esp$, podemos
demostrar para el espacio $D$ de infinitesimales tanto que $\forall x\in D, \neg\neg(x=0)$ como que $\neg(\forall x\in D, x=0)$,
por lo que $\esp$ es bivaluado, pero no es booleano. Podemos evidenciar que bivaluado no implica booleano, 
así como booleano no implica bivaluado, con ejemplos más sencillos y abundantes.

\begin{proposition}
  El topos $\con\times  \con$ es booleano y no bivaluado. 
\end{proposition}
\begin{proof}
  En este topos los objetos son pares de conjuntos $(X,Y)$, y los morfismos son pares de funciones 
  $(f,g)$, donde $f\colon X\to X'$ y $g\colon Y\to Y'$. Es sencillo comprobar que $(1,1)$ es el objeto terminal y que
  el clasificador de subobjetos es $(2,2)$, por lo que $\sub(1,1)\cong \homs((1,1),(2,2))\cong 4$ y el topos no es bivaluado,
  hay cuatro valores de verdad. Por otro lado, el coproducto se calcula como el
  coproducto de $\con$
  entrada a entrada,
  así que, $(1,1)+(1,1)=(1+1,1+1)=(2,2)$ y el topos es booleano. Para ver que el axioma~\ref{ax:gds6} no se cumple en este topos, 
  basta notar que el objeto $(1,0)$ no es el objeto inicial, sin embargo, no tiene puntos.
\end{proof}

\begin{proposition}
  Si $M$ es una categoría de un solo objeto y una flecha $f$ distinta a la identidad tal que 
  $f\circ f = f$, el topos $\con^M$ es bivaluado y no booleano.
\end{proposition}
\begin{proof}
  Los objetos de $\con^M$ son pares $(X,h)$ donde $X$ es un conjunto y $h\colon X\to X$ es 
  una función idempotente. Un morfismo $(X,h)\to (Y,k)$ es una función $g\colon X\to Y$ tal que 
  $g\circ h = k \circ g$. El objeto terminal es $(1,\id_1)$ y, al tratarse de un topos de pregavillas de $M$ un monoide, el 
  clasificador de subobjetos es el conjunto de ideales $\{\varnothing,\{f\},M\}\cong 3$, dotado de una endofunción $f_{\Omega}$ que manda a $0$ a $0$, $1$ a $2$ y $2$ a $2$. 
  Como el objeto terminal es un bucle, $(3,f_{\Omega})$ solo tiene dos elementos globales, así que el topos es bivaluado. 
  Ahora, en un topos de pregavillas los colímites se calculan puntualmente, de manera que $(1,\id_1)+(1,\id_1)=(2,\id_2)$ y
  por lo tanto, el topos no es booleano.
\end{proof}

\begin{theorem}
  En un topos cualquiera, si $A\in \sub(1)$ y $A$ tiene un punto $a\colon 1\to A$, entonces $A \cong 1$.
\end{theorem}
\begin{proof}
  Sea $A$ un subobjeto de $1$, con $m\colon A\to 1$ monomorfismo. Juntando nuestras hipótesis obtenemos al siguiente diagrama:
  \begin{equation*}
    \begin{tikzcd}
      1\ar{r}{a}\ar[swap]{rrd}{\text{id}} & A\ar{r}{m} & 1\ar[shift left = 0.3ex, no head]{d}{}\ar[shift right = 0.3ex, no head]{d}{} \\
       & {} & 1.
    \end{tikzcd}
  \end{equation*}
  El diagrama conmuta pues acaba en el terminal y como la composición es igual a la identidad
  en $1$, $m$ es un epimorfismo que se escinde aparte de ser mono. Por tanto, $m$ es un isomorfismo y $A\cong 1$.
\end{proof}

\begin{proposition}
  \pdfcomment{En un topos arbitrario, el Nullstellensatz es equivalente a que
  el topos sea bivaluado y los soportes se escindan.}
    Nullstellensatz si y solo si bivaluado y los soportes se escinden.
\end{proposition}
\begin{proof}
    Como vimos anteriormente, Nullstellensatz implica que el topos es bivaluado \ref{obs:nsbivaluado}, entonces
    el soporte de un objeto o es $0$ o es $1$ y en ambos casos tenemos garantizada una flecha del soporte al objeto.
    Ahora supongamos que el topos es bivaluado y los soportes se escinden. Sea $A$ un objeto,
    $\supp(A)$ es $0$ o $1$, pero los soportes se escinden entonces o $A$ es inicial (en una 
    categoría cartesianamente cerrada, $A\to 0$ implica $A\cong 0$) o bien, $A$ tiene un elemento. 
\end{proof}

Con esta proposición nos damos cuenta que la diferencia entre los axiomas de Mclarty y ETCS es si 
se asume o no el axioma de Elección.

Para una categoría de espacios, en la que se pueda llevar acabo el cálculo infinitesimal, 
Mclarty propone los siguientes axiomas.

\begin{axioma}[label=ax:gds1]
Existe un espacio, $R$, que cuenta con una estructura de anillo conmutativo con unidad.
\end{axioma}
\begin{axioma}[label=ax:gds2]
El anillo $R$ no es trivial, es decir, $1\not = 0$.
\end{axioma}
\begin{axioma}[label=ax:gds3]
En la lógica interna de $R$, $R$ es un campo.
\end{axioma}
Y el axioma de Kock-Lawvere \cite{lawvere1979categorical}, para el cual es necesario definir
al espacio de infinitesimales $D$, el cual será el igualador de $0\colon R\to R$, el mapa constante $0$, y $(-)^2\colon R\to R$.
Equivalentemente, $D$ es el subespacio (subobjeto) de $R$, $D= \{d \in R\colon d^2=0 \}$.
\begin{axioma}[label=ax:gds4]
$(\forall f \in R^D)(\exists ! b \in R)(\forall d \in D) (f(d) = f(0)+d \cdot b)$.
\end{axioma}
Intuitivamente, \ref{ax:gds4} dice que toda función de $D$ a $R$ es lineal. A continuación 
mostramos algunas propiedades de los infinitesimales y que estos cuatro axiomas ya nos 
permiten hacer cálculo básico.

\begin{definition}[label=def:discreto]
  Un espacio $E$ es discreto si y solo si es decidible, es decir, $(x=y) \lor \lnot(x = y)$ es verdad con $x,y$ variables
  variando sobre $E$.
\end{definition}

Mclarty le pide más axiomas a Espacios.

\begin{axioma}[label=ax:gds5]
Existe un espacio $N$ que es un objeto de números naturales.
\end{axioma}

Se puede demostrar que $N$ es es un espacio discreto.

\begin{definition}[label=puntos]
Los puntos de un espacio $E$ son sus elementos globales, morfismos de la forma $1\to E$
\end{definition}

\begin{axioma}[label=ax:gds6]
Para todo espacio $E$, o $E \cong  0 $ o bien $E$ tiene al menos un punto.
\end{axioma}

Esto implica que si $b_1, b_2\in B$ entonces $(b_1=b_2)\lor \lnot(b_1=b_2)$ es verdad incluso si $B$ no es discreto.

\begin{axioma}[label=ax:gds7]
Todo espacio $E$ tiene un único (salvo isomorfismo) subespacio $c_E\colon\Gamma E\to E$ discreto tal que todo punto de $E$ está en $\Gamma E$. 
\end{axioma} 